\section*{Bab \ref{ch:b2:comparison} --- Perbandingan Fungsi}

\begin{solution}{exo:b2:comparison:1}
$\sqrt{x^2 + x + 1} = x\sqrt{1 + \tfrac1x + \tfrac{1}{x^2}} = x +
\frac12 + \frac38\cdot\frac1x + o\bigl(\frac1x\bigr)$ (lewat uraian
binomial: $\frac12 u - \frac18 u^2$ dengan $u = \frac1x +
\frac{1}{x^2}$ memberi $\frac{1}{2x} + \frac{1}{2x^2} -
\frac{1}{8x^2} = \frac{1}{2x} + \frac{3}{8x^2}$, lalu dikalikan
$x$).

$\ln(x^2 + x) - 2\ln x = \ln\bigl(1 + \tfrac1x\bigr) = \frac1x -
\frac{1}{2x^2} + o\bigl(\frac{1}{x^2}\bigr)$.

Untuk fungsi ketiganya: uraikan tiap faktornya,
\[
\frac{x + \sin x}{x - \ln x}
= \Bigl(1 + \frac{\sin x}{x}\Bigr)
\Bigl(1 + \frac{\ln x}{x} + \frac{(\ln x)^2}{x^2} +
O\Bigl(\frac{(\ln x)^3}{x^3}\Bigr)\Bigr).
\]
Urutkan sumbangannya pada skala di $+\infty$: $\frac{\ln x}{x}
\gg \frac{1}{x} \geq \bigl|\frac{\sin x}{x}\bigr| \gg \frac{(\ln
x)^2}{x^2}$. Jadi dua suku sesudah suku utama $1$ adalah
$\frac{\ln x}{x}$, lalu suku berayun terbatas $\frac{\sin
x}{x}$:
\[
\frac{x + \sin x}{x - \ln x}
= 1 + \frac{\ln x}{x} + \frac{\sin x}{x}
+ O\Bigl(\frac{(\ln x)^2}{x^2}\Bigr).
\]
\end{solution}

\begin{solution}{exo:b2:comparison:2}
Ambil $f(t) = t^\alpha$ ($t \geq 1$).

$\alpha > -1$: divergen, dan menurut
\cref{thm:b2:comparison:seriesintegral}~(2), $\sum_{k\leq n}
k^\alpha = \frac{n^{\alpha+1}}{\alpha+1} + C + o(1)$ bila $\alpha <
0$ (yakni saat $f$ turun); sedangkan untuk $\alpha \geq 0$ (saat $f$
naik) pengapitan yang sama dengan ketaksamaan terbalik memberi
$\sum_{k \leq n} k^\alpha \sim \frac{n^{\alpha + 1}}{\alpha + 1}$.

$\alpha = -1$: $H_n = \ln n + \gamma + o(1)$
(\cref{ex:b2:comparison:harmonic}).

$\alpha < -1$: konvergen, dengan sisa $\sum_{k > n} k^\alpha
\sim \frac{n^{\alpha+1}}{-(\alpha+1)}$ menurut pengapitan (1) (sebab
kedua batas integralnya setara dengan nilai itu).
\end{solution}

\begin{solution}{exo:b2:comparison:3}
Misalkan $v_n = H_n - \ln n - \gamma \to 0$. Maka
\[
v_n - v_{n+1} = \ln\frac{n+1}{n} - \frac{1}{n+1}
= \Bigl(\frac1n - \frac{1}{2n^2}\Bigr) - \Bigl(\frac1n -
\frac{1}{n^2}\Bigr) + O\Bigl(\frac{1}{n^3}\Bigr)
= \frac{1}{2n^2} + O\Bigl(\frac{1}{n^3}\Bigr),
\]
dengan memakai $\frac{1}{n+1} = \frac1n - \frac{1}{n^2} + O(n^{-3})$.
Karena $v_n \to 0$, teleskopkan ekornya:
\[
v_n = \sum_{k \geq n} (v_k - v_{k+1})
= \sum_{k\geq n} \Bigl(\frac{1}{2k^2} + O(k^{-3})\Bigr)
= \frac{1}{2n} + O\Bigl(\frac{1}{n^2}\Bigr),
\]
menurut \cref{thm:b2:comparison:seriesintegral}~(1) yang diterapkan pada
$t^{-2}$ (sisanya $\sim \frac1n$, lalu diparuhkan) dan pada $t^{-3}$.
Karenanya $H_n = \ln n + \gamma + \frac{1}{2n} + o(\frac1n)$.
\end{solution}

\begin{solution}{exo:b2:comparison:4}
Stirling tiga kali:
\[
\frac{(3n)!}{(n!)^3}
\sim \frac{\sqrt{6\pi n}\,(3n/\eu)^{3n}}
{(2\pi n)^{3/2}\,(n/\eu)^{3n}}
= \frac{\sqrt{6}\; 27^{\,n}}{2\pi n} \cdot
\frac{1}{\sqrt{2\pi n}}\cdot\sqrt{2\pi n}\;
= \frac{\sqrt3\,27^n}{2\pi n} .
\]
(Dengan cermat: $\frac{\sqrt{6\pi n}}{(2\pi n)^{3/2}} =
\frac{\sqrt6}{(2\pi n)\sqrt{2\pi n}}\sqrt{\pi n} =
\frac{\sqrt3}{2\pi n}$.)

$\dfrac{n!}{n^n} \sim \sqrt{2\pi n}\,\eu^{-n}$.

Dan $\sqrt[n]{n!} = \exp\bigl(\frac{\ln n!}{n}\bigr)$ dengan $\ln n! =
n\ln n - n + \frac12\ln(2\pi n) + o(1)$:
\[
\sqrt[n]{n!} = \exp\Bigl(\ln n - 1 + \frac{\ln(2\pi n)}{2n} +
o\Bigl(\frac{\ln n}{n}\Bigr)\Bigr)
= \frac{n}{\eu}\Bigl(1 + \frac{\ln(2\pi n)}{2n} +
o\Bigl(\frac{\ln n}{n}\Bigr)\Bigr).
\]
\end{solution}

\begin{solution}{exo:b2:comparison:5}
Fungsi $g(x) = x^n + x - 1$ naik tegas pada $\intcc{0}{1}$ dari $-1$
ke $1$: jadi akarnya tunggal, yakni $x_n$. Karena $x_n^n = 1 - x_n \in
\intoo{0}{1}$: seandainya $x_n \leq c < 1$ sepanjang suatu barisan
bagian, maka $x_n^n \leq c^n \to 0$, sehingga $1 - x_n \to 0$:
kontradiksi dengan $x_n \leq c$. Karenanya $x_n \to 1$.

Tulis $x_n = 1 - \varepsilon_n$ dengan $\varepsilon_n \to 0^+$.
Persamaannya berbunyi $(1 - \varepsilon_n)^n = \varepsilon_n$, yakni
\[
n\ln(1 - \varepsilon_n) = \ln \varepsilon_n
\quad\Longrightarrow\quad
-n\varepsilon_n\bigl(1 + o(1)\bigr) = \ln\varepsilon_n .
\]
Jadi $n\varepsilon_n = -\ln\varepsilon_n\,(1 + o(1)) \to +\infty$, dan
setelah logaritmanya diambil lagi: $\ln n + \ln\varepsilon_n =
\ln(-\ln\varepsilon_n) + o(1)$. Karena $\ln(-\ln \varepsilon_n) =
o(\ln(1/\varepsilon_n))$, hal ini memberi $\ln\varepsilon_n \sim -\ln
n$, sehingga $\varepsilon_n = \frac{-\ln\varepsilon_n}{n}(1 + o(1))
\sim \frac{\ln n}{n}$:
\[
x_n = 1 - \frac{\ln n}{n} + o\Bigl(\frac{\ln n}{n}\Bigr) .
\]
\end{solution}

\begin{solution}{exo:b2:comparison:6}
Relasi eksaknya: $\cot y_n = x_n = n\pi + \frac\pi2 - y_n$, dengan $y_n
\sim \frac{1}{n\pi}$ (\cref{ex:b2:comparison:tan}). Uraikan $\cot y
= \frac1y - \frac y3 + O(y^3)$:
\[
\frac{1}{y_n} - \frac{y_n}{3} + O(y_n^3) = n\pi + \frac\pi2 - y_n
\quad\Longrightarrow\quad
\frac{1}{y_n} = n\pi + \frac\pi2 + O\Bigl(\frac1n\Bigr),
\]
(sebab suku $-y_n$ dan $-\frac{y_n}{3}$ bernilai $O(\frac1n)$). Balikkan:
\[
y_n = \frac{1}{n\pi}\cdot\frac{1}{1 + \frac{1}{2n} + O(n^{-2})}
= \frac{1}{n\pi}\Bigl(1 - \frac{1}{2n} + O\Bigl(\frac{1}{n^2}\Bigr)\Bigr)
= \frac{1}{n\pi} - \frac{1}{2n^2\pi} + O\Bigl(\frac{1}{n^3}\Bigr).
\]
Karenanya
\[
x_n = n\pi + \frac{\pi}{2} - y_n
= n\pi + \frac\pi2 - \frac{1}{n\pi} + \frac{1}{2n^2\pi} +
o\Bigl(\frac{1}{n^2}\Bigr).
\]
\end{solution}

\begin{solution}{exo:b2:comparison:7}
Pisahkan dua suku terbesarnya:
\[
\sum_{k=0}^{n} k! = n! + (n-1)! + \sum_{k \leq n-2} k! ,
\qquad
\sum_{k\leq n-2} k! \leq (n-1)\,(n-2)! = (n-1)! .
\]
Jadi $1 \leq \frac{1}{n!}\sum k! \leq 1 + \frac{2}{n}$: sehingga
limitnya $1$. Diperhalus: $\frac{(n-1)!}{n!} = \frac1n$ dan batas kasar
$\sum_{k \leq n-2}k! \leq (n-1)!$ dapat dipertajam dengan cara yang sama:
$\sum_{k\leq n-2} k! = (n-2)!\,(1 + O(\frac1n)) = O\bigl(\frac{n!}{n^2}\bigr)$.
Karenanya
\[
\sum_{k=0}^{n} k! = n!\Bigl(1 + \frac1n + O\Bigl(\frac{1}{n^2}\Bigr)\Bigr).
\]
\end{solution}

\begin{solution}{exo:b2:comparison:8}
Barisan $(u_n)$ naik; seandainya ia terbatas, ia akan konvergen ke
$\ell$ dengan $\ell = \ell + \frac1\ell$: mustahil. Jadi $u_n \to \infty$.

Kuadratnya: $u_{n+1}^2 = u_n^2 + 2 + u_n^{-2}$, sehingga
\[
u_n^2 = u_0^2 + 2n + \sum_{k=0}^{n-1} \frac{1}{u_k^2} .
\]
Jumlahnya bernilai $o(n)$ (karena sukunya menuju $0$, lewat Cesàro),
jadi $u_n^2 \sim 2n$ dan $u_n \sim \sqrt{2n}$.

Perhalusannya: $\frac{1}{u_k^2} \sim \frac{1}{2k}$, jadi lewat
perbandingan (\cref{thm:b2:comparison:seriesintegral}, atau lewat yang
setara bagi jumlah parsial deret positif) berlaku $\sum_{k<n} u_k^{-2}
\sim \frac12 \ln n$. Karenanya
\[
u_n^2 = 2n + \frac{\ln n}{2}\,(1 + o(1)) + O(1)
\quad\Longrightarrow\quad
u_n = \sqrt{2n}\sqrt{1 + \frac{\ln n}{4n} + o\Bigl(\frac{\ln
n}{n}\Bigr)}
= \sqrt{2n}\Bigl(1 + \frac{\ln n}{8n} + o\Bigl(\frac{\ln
n}{n}\Bigr)\Bigr).
\]
\end{solution}

\begin{solution}{exo:b2:comparison:9}
Keluarkan faktor $n$ lalu tetapkan $t = \ln n$:
\[
S_n = \frac1n \sum_{k=1}^{n} \frac{1}{1 + t\,\frac kn} .
\]
Untuk $t$ yang tetap, jumlahnya adalah jumlah Riemann bagi $u \mapsto
\frac{1}{1 + tu}$ pada $\intcc{0}{1}$; fungsinya monoton pada $u$, jadi
jumlah Riemann itu terapit oleh integralnya yang digeser satu petak:
\[
\int_0^1 \frac{\dd u}{1 + tu} - \frac1n
\leq S_n \leq \int_0^1 \frac{\dd u}{1 + tu} + \frac1n
\]
(yakni perbandingan jumlah Riemann sebuah fungsi monoton dengan
integralnya, yang sahih untuk tiap $n$ dengan $t = \ln n$-nya sendiri).
Kini $\int_0^1 \frac{\dd u}{1 + tu} = \frac{\ln(1 + t)}{t}$, dan
$\frac1n = o\bigl(\frac{\ln t}{t}\bigr)$: karenanya
\[
S_n = \frac{\ln(1 + \ln n)}{\ln n} + O\Bigl(\frac 1n\Bigr)
\;\sim\; \frac{\ln\ln n}{\ln n} .
\]
\end{solution}

\begin{solution}{exo:b2:comparison:10}
Kesamaannya: $(\ln n)^{\ln n} = \eu^{\ln n\,\ln\ln n} =
\bigl(\eu^{\ln n}\bigr)^{\ln\ln n} = n^{\ln\ln n}$. Pengurutannya:
bandingkan logaritmanya. $\ln(n^2) = 2\ln n$; $\ln\bigl((\ln
n)^{\ln n}\bigr) = \ln n\ln\ln n$; $\ln(2^n) = n\ln2$;
$\ln(n!) = n\ln n - n + O(\ln n)$ (lewat Stirling, atau lewat
pengapitan yang lebih kasar $\ln n! \sim n\ln n$); $\ln(n^n) = n\ln n$.
Karena $2\ln n = o(\ln n\ln\ln n)$, $\ln n\ln\ln n = o(n)$, $n\ln 2 =
o(n\ln n - n)$, dan $n \ln n - n \sim n\ln n$ padahal $n! / n^n \to
0$ (sebab selisih logaritmanya $-n + O(\ln n) \to -\infty$), diperoleh
\[
n^2 = o\bigl((\ln n)^{\ln n}\bigr),\quad
(\ln n)^{\ln n} = o(2^n),\quad
2^n = o(n!),\quad
n! = o(n^n).
\]
(Untuk tiap langkahnya: selisih logaritmanya menuju
$+\infty$, jadi rasionya menuju $0$.)
\end{solution}

\begin{solution}{exo:b2:comparison:11}
Uraikan $\frac1{k^2} = \frac1{k(k+1)} + \frac1{k^2(k+1)}$ lalu
jumlahkan untuk $k > n$:
\[
\sum_{k>n}\frac1{k^2} = \frac1{n+1} +
\sum_{k>n}\frac{1}{k^2(k+1)} ,
\]
dengan jumlah pertamanya berteleskop secara eksak (sebab $\frac1{k(k+1)}
= \frac1k - \frac1{k+1}$). Untuk yang kedua: $\frac{1}{k^2(k+1)} =
\frac1{k^3} + O\bigl(\frac1{k^4}\bigr)$ (sebab $\frac{1}{k^2(k+1)} -
\frac1{k^3} = \frac{-1}{k^3(k+1)}$), dan menurut
perbandingan integral $\sum_{k>n}\frac1{k^3} = \frac1{2n^2} +
O\bigl(\frac1{n^3}\bigr)$ serta $\sum_{k>n}\frac1{k^4} =
O\bigl(\frac1{n^3}\bigr)$. Karenanya
\[
\sum_{k>n}\frac1{k^2}
= \frac1{n+1} + \frac{1}{2n^2} + O\Bigl(\frac1{n^3}\Bigr)
= \frac1n - \frac1{n^2} + \frac{1}{2n^2} +
O\Bigl(\frac1{n^3}\Bigr)
= \frac1n - \frac{1}{2n^2} + O\Bigl(\frac1{n^3}\Bigr),
\]
dengan memakai $\frac1{n+1} = \frac1n - \frac1{n^2} +
O\bigl(\frac1{n^3}\bigr)$.
\end{solution}

\begin{solution}{exo:b2:comparison:12}
Barisan $(u_n)$ naik; seandainya ia terbatas, ia akan konvergen ke
$\ell$ yang hingga dengan $\ell = \ell + \eu^{-\ell}$: mustahil. Jadi
$u_n \to \infty$. Misalkan $v_n = \eu^{u_n} \to \infty$: maka
\[
v_{n+1} = \eu^{u_n + \eu^{-u_n}} = v_n\,\eu^{1/v_n}
= v_n\Bigl(1 + \frac1{v_n} + \frac1{2v_n^2} +
O\bigl(v_n^{-3}\bigr)\Bigr)
= v_n + 1 + \frac{1}{2v_n} + O\bigl(v_n^{-2}\bigr).
\]
Menjumlahkan $v_{k+1} - v_k = 1 + O(1)$ lebih dulu memberi $v_n = n +
O(n)$, jadi akhirnya $v_n \geq cn$; lalu menjumlahkannya ulang dengan
$\frac1{2v_k} = O(\frac1k)$ memberi $v_n = n + O(\ln n)$. Satu lintasan
lagi: $\frac{1}{2v_k} = \frac{1}{2k}\bigl(1 +
O\bigl(\tfrac{\ln k}k\bigr)\bigr)$, sehingga
\[
v_n = n + \sum_{k<n}\frac1{2k} + O(1) = n + \frac{\ln n}2 +
O(1).
\]
Akhirnya $u_n = \ln v_n = \ln n + \ln\Bigl(1 + \frac{\ln n}{2n} +
O\bigl(\tfrac1n\bigr)\Bigr) = \ln n + \frac{\ln n}{2n} +
O\bigl(\tfrac1n\bigr)$.
\end{solution}

\begin{solution}{pb:b2:comparison:1}
\textbf{1.} Setelah kedua uraiannya dikurangkan: $\sum_i (c_i -
c_i')\varphi_i = o(\varphi_k)$. Jika ada koefisien yang berbeda, misalkan
$i_0$ yang pertama: membaginya dengan $\varphi_{i_0}$ lalu memakai
$\varphi_j = o(\varphi_{i_0})$ untuk $j > i_0$ memberi $c_{i_0} -
c_{i_0}' = o(1)$: jadi nol, kontradiksi. Untuk uraiannya: dengan
$u = \frac{\ln x}x \to 0$,
\[
\frac{1}{x - \ln x} = \frac1x\cdot\frac{1}{1 - u}
= \frac1x\bigl(1 + u + u^2 + O(u^3)\bigr)
= \frac1x + \frac{\ln x}{x^2} + \frac{(\ln x)^2}{x^3} +
o\Bigl(\frac{(\ln x)^2}{x^3}\Bigr).
\]
Tak ada suku $\frac c{x^2}$ yang muncul karena uraiannya berupa
deret geometri dalam $u = \frac{\ln x}{x}$: tiap sukunya mengusung
pangkat $\ln x$ sebanyak pangkat $\frac1x$ melampaui yang pertama;
jadi anak tangga $\frac1{x^2}$ (yakni koefisien $(\ln x)^0$) memang
tak ada, dengan koefisien $0$.

\textbf{2.} Fungsi $f(x) = \eu^x + x$ bersifat \omterm{def:b2:metric:continuity}{kontinu}, naik
tegas, dengan limit $-\infty$ dan $+\infty$: jadi ia bijeksi
$\R \to \R$, sehingga $x_n = f^{-1}(n)$ ada dan tunggal, dan
$x_n \to +\infty$ (sebab $f^{-1}$ naik ke $+\infty$). Dari
$\eu^{x_n} = n - x_n$: $x_n = \ln(n - x_n) \leq \ln n$, jadi
$x_n/n \to 0$ dan $x_n = \ln n + \ln(1 - x_n/n) = \ln n + o(1)
\sim \ln n$.

\textbf{3.} Tulis $u_n = x_n/n$. Lintasan kedua: $u_n = \frac{\ln n
+ o(1)}{n}$, sehingga
\[
x_n = \ln n + \ln(1 - u_n) = \ln n - u_n + O(u_n^2)
= \ln n - \frac{\ln n}{n} + o\Bigl(\frac{\ln n}n\Bigr).
\]
Lintasan ketiga: kini $u_n = \frac{\ln n}{n} - \frac{\ln n}{n^2} +
o\bigl(\frac{\ln n}{n^2}\bigr)$, dan $\ln(1 - u_n) = -u_n -
\frac{u_n^2}2 + O(u_n^3)$:
\[
x_n = \ln n - \frac{\ln n}n + \frac{\ln n}{n^2}
- \frac{(\ln n)^2}{2n^2} + o\Bigl(\frac{(\ln
n)^2}{n^2}\Bigr)
= \ln n - \frac{\ln n}{n} - \frac{(\ln n)^2}{2n^2} +
o\Bigl(\frac{(\ln n)^2}{n^2}\Bigr),
\]
dengan suku $\frac{\ln n}{n^2}$ terserap ke dalam
$o\bigl(\frac{(\ln n)^2}{n^2}\bigr)$.

\textbf{4.} Di $n = 1000$: $\ln 1000 \approx 6.90776$ (galatnya
$7\cdot10^{-3}$); dua suku memberi $6.90085$ (galatnya
$2\cdot10^{-5}$); tiga suku memberi $6.90082$ (galatnya di bawah
$10^{-5}$), terhadap $x_{1000} \approx 6.90083$. Jadi tiap lintasan
membeli kira-kira faktor $\frac{\ln n}{n}$ yang diramalkan.

\textbf{5.} Dua kali pengintegralan parsial, berawal dari kanan:
dengan $\frac{\dd}{\dd t}\bigl[\tfrac12t(1-t)\bigr] = \tfrac12 -
t$ dan $t(1-t)$ yang nol di kedua ujungnya,
\[
\frac12\int_0^1 t(1-t)g''(t)\dd t
= -\int_0^1\Bigl(\frac12 - t\Bigr)g'(t)\dd t
= -\Bigl[\Bigl(\frac12 - t\Bigr)g\Bigr]_0^1 - \int_0^1 g
= \frac{g(0) + g(1)}2 - \int_0^1 g .
\]
Setelah ditata ulang, inilah kesamaan yang dinyatakan tadi.

\textbf{6.} Hitunglah pertambahannya, lalu terapkan pertanyaan 5 pada
$g(t) = f(n + t)$:
\begin{align*}
E_{n+1} - E_n
&= f(n{+}1) - \int_n^{n+1}\!f - \frac{f(n{+}1) - f(n)}2 \\
&= \frac{f(n) + f(n{+}1)}2 - \int_n^{n+1}\!f
= \frac12\int_0^1 t(1-t)f''(n+t)\dd t .
\end{align*}
Karena $0 \leq t(1-t) \leq \frac14$: berlaku $\abs{E_{n+1} - E_n} \leq
\frac18\int_n^{n+1}\abs{f''}$, yang jumlahnya atas $n$ konvergen menurut
hipotesisnya: jadi $(E_n)$ konvergen (sebab pertambahannya terjumlahkan
secara mutlak) ke suatu $E$, dengan
\[
\abs{E - E_n} \leq \sum_{k\geq n}\abs{E_{k+1} - E_k} \leq
\frac18\int_n^\infty\abs{f''} .
\]

\textbf{7.} Untuk $f(t) = \frac1t$: $f''(t) = \frac2{t^3}$ dan
$\int_1^\infty\abs{f''} = 1 < \infty$. Menurut pertanyaan 6:
\[
H_n = \ln n + \frac{1 + \frac1n}{2} + E + (E_n - E)
= \ln n + \Bigl(E + \frac12\Bigr) + \frac1{2n} +
\varepsilon_n,
\]
dengan $\abs{\varepsilon_n} = \abs{E_n - E} \leq
\frac18\int_n^\infty\frac{2\dd t}{t^3} = \frac1{8n^2}$.
Membandingkannya dengan $H_n = \ln n + \gamma + o(1)$
(\cref{ex:b2:comparison:harmonic}) mengenali $E + \frac12 =
\gamma$.

\textbf{8.} Dari rumus pertambahan pada pertanyaan 6,
\[
\varepsilon_n = E_n - E = -\sum_{k\geq n}\frac12\int_0^1
t(1-t)\,\frac{2\,\dd t}{(k+t)^3}
= -\sum_{k \geq n}\Bigl(\frac1{k^3}\int_0^1t(1-t)\dd t +
O\Bigl(\frac1{k^4}\Bigr)\Bigr),
\]
dengan memakai $\frac{1}{(k+t)^3} = \frac1{k^3} +
O\bigl(\frac1{k^4}\bigr)$ secara seragam untuk $t \in \intcc01$. Dengan
$\int_0^1 t(1-t) = \frac16$ dan
$\sum_{k\geq n}\frac1{k^3} \sim \frac{1}{2n^2}$
(\cref{thm:b2:comparison:seriesintegral}):
\[
\varepsilon_n = -\frac16\cdot\frac{1}{2n^2} +
o\Bigl(\frac1{n^2}\Bigr) = -\frac{1}{12n^2} +
o\Bigl(\frac{1}{n^2}\Bigr).
\]

\textbf{9.} Untuk $f = \ln$: $f''(t) = -\frac1{t^2}$, yang
terintegralkan secara mutlak. Pertanyaan 6 memberi
\[
\ln n! = \int_1^n\ln t\,\dd t + \frac{\ln n}2 + E + O\Bigl(
\frac1{8}\int_n^\infty\frac{\dd t}{t^2}\Bigr)
= \Bigl(n + \frac12\Bigr)\ln n - n + 1 + E +
O\Bigl(\frac1n\Bigr),
\]
jadi $d_n = 1 + E + O\bigl(\frac1n\bigr)$: yakni kekonvergenan $(d_n)$
--- Langkah 1 pada \cref{thm:b2:comparison:stirling} --- ditambah
lajunya $O(1/n)$. (Nilai limit menurut Stirling memberi $E =
\ln\sqrt{2\pi} - 1$.)

\textbf{10.} Untuk $f(t) = t^{-1/2}$: $f''(t) = \frac34 t^{-5/2}$, yang
terintegralkan secara mutlak. Pertanyaan 6 memberi
\[
\sum_{k=1}^n \frac1{\sqrt k}
= 2\sqrt n - 2 + \frac{1 + \frac1{\sqrt n}}2 + E +
O\bigl(n^{-3/2}\bigr)
= 2\sqrt n + c + \frac{1}{2\sqrt n} +
O\bigl(n^{-3/2}\bigr),
\]
dengan $c = E - \frac32$. Di $n = 10^4$: $2\sqrt n = 200$, $c
\approx -1.46035$, $\frac1{2\sqrt n} = 0.005$: jadi ramalannya
$198.54465$, dan memang $\sum_{k\leq10^4}k^{-1/2} =
198.544645\dots$ --- tiga suku, tujuh angka.

\textbf{11.} Pemetaan $t \mapsto t\ln t$ \omterm{def:b2:metric:continuity}{kontinu} dan naik
tegas pada $\intco1\infty$ (sebab turunannya $\ln t + 1 \geq 1$),
dari $0$ ke $+\infty$: jadi ada $x_n$ yang tunggal, dan $x_n \to
\infty$ (sebab jika tidak, $x_n\ln x_n$ akan tetap terbatas). Setelah
logaritmanya diambil pada $x_n\ln x_n = n$: $\ln x_n + \ln\ln x_n = \ln
n$; dan karena $\ln\ln x_n = o(\ln x_n)$, membaginya dengan $\ln x_n$
memberi $\frac{\ln n}{\ln x_n} \to 1$: jadi $\ln x_n \sim \ln n$.

\textbf{12.} Dari $x_n = \frac{n}{\ln x_n}$ dan $\ln x_n \sim
\ln n$: $x_n \sim \frac{n}{\ln n}$. Lintasan berikutnya: $\ln\ln x_n =
\ln\bigl(\ln n\,(1 + o(1))\bigr) = \ln\ln n + o(1)$, jadi $\ln
x_n = \ln n - \ln\ln n + o(1)$ dan
\[
x_n = \frac{n}{\ln n - \ln\ln n + o(1)}
= \frac{n}{\ln n}\cdot\frac{1}{1 - \frac{\ln\ln n +
o(1)}{\ln n}}
= \frac{n}{\ln n}\Bigl(1 + \frac{\ln\ln n}{\ln n} +
o\Bigl(\frac{\ln\ln n}{\ln n}\Bigr)\Bigr).
\]

\textbf{13.} Di $n = 10^6$: $\frac{n}{\ln n} \approx 72\,382$
(meleset $18\%$), dua suku memberi $\approx 86\,140$ (meleset
$1.9\%$), terhadap nilai sejatinya $x \approx 87\,848$. Keuntungan tiap
lintasannya hanya sebesar faktor $\frac{\ln\ln n}{\ln n} \approx
\frac{2.63}{13.8} \approx 0.19$: jadi skala logaritmik konvergen
dengan kelambanan yang menjengkelkan --- kenyataan hidup di mana pun
bilangan prima terlibat.

\textbf{14.} Ada tepat $n$ bilangan prima $\leq p_n$ (yakni
$p_1, \dots, p_n$): jadi $\pi(p_n) = n$. Teorema bilangan prima
(yang diterima tanpa bukti; jilid Tahun ke-3) memberi $n = \pi(p_n) \sim
\frac{p_n}{\ln p_n}$, yakni $p_n \sim n\ln p_n$: dan inilah
persamaan $x\ln x \approx n$ yang dibaca terbalik. Setelah logaritmanya
diambil: $\ln p_n = \ln n + \ln\ln p_n + o(1)$, dan $\ln\ln p_n =
o(\ln p_n)$ memaksa $\ln p_n \sim \ln n$ seperti pada pertanyaan 11.
Setelah disulihkan kembali:
\[
p_n \sim n\ln p_n = n\,\ln n\,\frac{\ln p_n}{\ln n} \sim n\ln
n .
\]

\textbf{15.} (a) Tetapkan $\varepsilon > 0$; untuk $k$ yang besar, $(1 -
\varepsilon)k\ln k \leq p_k \leq (1 + \varepsilon)k\ln k$. Lewat
perbandingan dengan $t\ln t$ yang naik
(dengan pengapitan bergaya \cref{thm:b2:comparison:seriesintegral}),
$\sum_{k\leq n}k\ln k = \int_1^n t\ln t\,\dd t + O(n\ln n) =
\frac{n^2\ln n}2 - \frac{n^2}4 + O(n\ln n) \sim
\frac{n^2\ln n}2$. Karenanya $\sum_{k\leq n}p_k = \frac{n^2\ln
n}{2}(1 + O(\varepsilon) + o(1))$ untuk setiap $\varepsilon$:
jadi $\sum_{k\leq n}p_k \sim \frac{n^2\ln n}2$. (b) Menurut teorema
bilangan prima, di antara bilangan bulat sampai $10^{100}$ ada
proporsi $\sim \frac{1}{\ln 10^{100}} = \frac1{230.26\dots}$
yang prima: jadi bilangan bulat acak seragam berangka $100$ bersifat
prima dengan peluang kira-kira $\frac1{230}$.

\textbf{16.} Pada $\intoo{n\pi}{n\pi + \frac\pi2}$, fungsi $g(x) = \tan
x - \frac1x$ \omterm{def:b2:metric:continuity}{kontinu} dan naik tegas (sebab $g' = 1 +
\tan^2x + \frac1{x^2} > 0$), dengan $g \to -\frac1{n\pi} < 0$ di
ujung kirinya dan $g \to +\infty$ di ujung kanannya: jadi ada tepat satu
akar $x_n$. Karena $\tan z_n = \tan x_n = \frac1{x_n} \to 0$
dengan $z_n \in \intoo{0}{\frac\pi2}$: maka $z_n \to 0^+$.

\textbf{17.} Kita punya $\tan z_n \sim z_n$ dan $\frac1{x_n} \sim
\frac1{n\pi}$: jadi $z_n \sim \frac1{n\pi}$.

\textbf{18.} Di sini $z_n = \arctan\frac1{x_n}$ dan $\arctan u = u +
O(u^3)$. Dengan $z_n = O(\frac1n)$:
\[
\frac1{x_n} = \frac{1}{n\pi}\cdot\frac1{1 + \frac{z_n}{n\pi}}
= \frac1{n\pi} - \frac{z_n}{n^2\pi^2} +
O\Bigl(\frac1{n^4}\Bigr)
= \frac1{n\pi} + O\Bigl(\frac1{n^3}\Bigr),
\]
jadi $z_n = \frac1{n\pi} + O\bigl(\frac1{n^3}\bigr)$: yakni anak tangga
$\frac{c}{n^2}$ berkoefisien $0$, sebab koreksi pertama pada
$\frac1{x_n}$ itu sendiri berukuran $\frac{z_n}{n^2} = O(n^{-3})$.

\textbf{19.} Masukkan $z_n = \frac1{n\pi} + O(n^{-3})$ ke dalam
ungkapan sebelumnya:
\[
\frac{1}{x_n} = \frac{1}{n\pi} - \frac{1}{n^3\pi^3} +
O\Bigl(\frac1{n^5}\Bigr),
\]
lalu $z_n = \arctan\frac1{x_n} = \frac1{x_n} -
\frac{1}{3}\Bigl(\frac1{x_n}\Bigr)^3 + O\Bigl(\frac1{n^5}\Bigr)
= \frac1{n\pi} - \frac{1}{n^3\pi^3} - \frac{1}{3n^3\pi^3} +
O\Bigl(\frac1{n^5}\Bigr)$:
\[
x_n = n\pi + \frac{1}{n\pi} - \frac{4}{3\pi^3n^3} +
O\Bigl(\frac1{n^5}\Bigr).
\]

\textbf{20.} Di $n = 3$: satu suku memberi $9.53088$, tiga suku
memberi $9.52929$, sedangkan akar sejatinya $9.52933$: jadi galatnya
$1.5\cdot10^{-3}$ dan $5\cdot10^{-5}$. Bandingkan: untuk $\tan x = x$
akarnya wajib membuat $\tan$ menjadi besar sekali, jadi ia memeluk ujung
\emph{kanan} $n\pi + \frac\pi2$ pada jendelanya, berjarak
$\sim\frac1{n\pi}$ sebelum asimtotnya; sedangkan untuk $x\tan x = 1$
akarnya wajib membuat $\tan$ menjadi sangat kecil, jadi ia duduk tepat
sesudah ujung \emph{kiri} $n\pi$, berjarak $\sim\frac1{n\pi}$ sesudah
nolnya. Metode yang sama, geografi yang bercermin.

\textbf{21.} Kita punya $\sin u < u$ pada $\intoo0\pi$ dan $\sin$
memetakan $\intoo0\pi$ ke dalam $\intoc01 \subseteq \intoo0\pi$: jadi
setelah satu langkah $u_1 \in \intoc{0}{1}$, lalu $(u_n)$ turun dan
terbatas di bawah oleh $0$, sehingga ia konvergen ke titik tetap $\sin$,
yakni ke $0$. Uraiannya: $\sin u = u(1 - \frac{u^2}6 +
o(u^2))$, sehingga
\[
\frac{1}{u_{n+1}^2} - \frac1{u_n^2}
= \frac{1}{u_n^2}\Bigl(\bigl(1 - \tfrac{u_n^2}6 +
o(u_n^2)\bigr)^{-2} - 1\Bigr)
= \frac{1}{u_n^2}\Bigl(\frac{u_n^2}{3} + o(u_n^2)\Bigr)
\longrightarrow \frac13 .
\]

\textbf{22.} Menurut Cesàro (jilid Tahun ke-1), rata-rata
pertambahannya konvergen ke limit yang sama:
\[
\frac{1}{n}\cdot\frac{1}{u_n^2}
= \frac1n\Bigl(\frac1{u_0^2} + \sum_{k=0}^{n-1}
\Bigl(\frac1{u_{k+1}^2} - \frac1{u_k^2}\Bigr)\Bigr)
\longrightarrow \frac13 ,
\]
jadi $u_n^2 \sim \frac3n$ dan, karena semua sukunya positif, $u_n
\sim \sqrt{3/n}$.

\textbf{23.} Menurut pertanyaan 7, $\abs{H_n - \ln n - \gamma -
\frac1{2n}} \leq \frac1{8n^2}$. Di $n = 10^6$ batas ini bernilai
$\frac{1}{8\cdot10^{12}} = 1.25\cdot10^{-13}$: jadi tiga suku yang
terhitung itu menyerahkan jumlah harmonik sejuta suku sampai tiga belas
angka, \omterm{def:b2:metric:complete}{lengkap} dengan sertifikat galat yang sepenuhnya cermat --- dan
itulah seluruh inti rumus asimtotik yang bersisa tersurat.

\textbf{24.} (a) Benar: $\ln u_n - \ln v_n = \ln\frac{u_n}{v_n}
\to 0$ sedangkan $\ln v_n \to +\infty$, jadi rasio logaritmanya
menuju $1$. (b) Salah: $u_n = n + 1 \sim v_n = n$, tetapi
$\eu^{u_n}/\eu^{v_n} = \eu \neq 1$. Kesetaraan memaklumi
galat aditif $o(1)$ pada eksponennya, bukan $O(1)$. (c) Salah:
$f(x) = x + \sin(x^2) \sim g(x) = x$ di $+\infty$, tetapi $f'(x)
= 1 + 2x\cos(x^2)$ berayun tanpa batas sedangkan $g' = 1$:
jadi turunan dua fungsi yang setara sama sekali tak harus sebanding.

\textbf{25.} Loop pada \cref{met:b2:comparison:implicit} berjalan
serupa tiga kali: setempatkan akarnya, sarikan suku
kasarnya, lalu umpankan kembali untuk orde berikutnya --- pada $\eu^x + x
= n$ (Bagian I), pada $x\ln x = n$ (Bagian III), dan pada $x\tan x = 1$
(Bagian IV). Koreksi trapesiumnya meningkatkan perbandingan deret
dengan integral dari ``selisihnya konvergen'' menjadi suku
$\frac{f(1) + f(n)}2$ yang tersurat dengan sisa $O(\int_n^\infty
\abs{f''})$ yang tersahkan --- yakni konstanta dan batang galat, bukan
sekadar kekonvergenan. Jembatan ke bilangan prima murni pembalikan:
teorema bilangan prima mengatakan $\pi(x)\ln x \approx x$, jadi $p_n$,
yang ditetapkan oleh $\pi(p_n) = n$, menyelesaikan sebuah persamaan
$x\ln x = n$ --- sehingga ia mewarisi asimtotiknya. Aturan (a) pada
pertanyaan 24 mengesahkan setiap peralihan dari $u_n \sim v_n$ ke $\ln
u_n \sim \ln v_n$ (pertanyaan 11 dan 14); sedangkan kesalahan (b) itulah
sebabnya kita tak pernah mengeksponensialkan kesetaraan. Puncaknya:
rumus Euler--Maclaurin sampai orde pertama (pertanyaan 6), dan hukum
asimtotik $p_n \sim n\ln n$ bagi bilangan prima ke-$n$ (pertanyaan 14).

\end{solution}
